\documentclass[a4paper,12pt]{article}
\usepackage[T2A]{fontenc}
\usepackage[utf8]{inputenc}
\usepackage[russian]{babel}
\usepackage{amsmath,amssymb,mathtools}
\usepackage{helvet}
\renewcommand{\familydefault}{\sfdefault}
\usepackage[margin=23mm,top=25mm,bottom=25mm]{geometry}
\usepackage{microtype,booktabs,array,tabularx,enumitem,xcolor,tikz,hyperref,fancyhdr}
\usetikzlibrary{calc}
\definecolor{accent}{HTML}{006E75}
\definecolor{ink}{HTML}{263238}
\definecolor{muted}{HTML}{657478}
\definecolor{wrong}{HTML}{B23B36}
\hypersetup{colorlinks=true,linkcolor=accent,pdfauthor={Лёвин Артём Александрович},pdftitle={Ревью домашней работы — Анастасия Павлова}}
\setlength{\parindent}{0pt}
\setlength{\parskip}{8pt}
\setlength{\headheight}{15pt}
\pagestyle{fancy}
\fancyhf{}
\fancyhead[L]{\small\color{muted}Математика · ревью домашней работы}
\fancyhead[R]{\small\color{muted}2 ноября 2026 г.}
\fancyfoot[C]{\footnotesize\color{muted}Лёвин Артём Александрович · эксклюзивно для Анастасии Павловой}
\fancyfoot[R]{\footnotesize\thepage}
\renewcommand{\headrulewidth}{0pt}
\renewcommand{\footrulewidth}{0pt}
\newcommand{\field}[1]{\vspace{3pt}{\color{accent}\bfseries #1}\par\vspace{-2pt}}
\newcommand{\errorlabel}{\textcolor{wrong}{\bfseries Ошибка:}}
\newcommand{\keylabel}{\textcolor{accent}{\bfseries Ключевая идея:}}
\newcommand{\resultlabel}{\textcolor{accent}{\bfseries Итог:}}
\begin{document}
\thispagestyle{empty}
\vspace*{29mm}
\begin{center}
{\footnotesize\color{muted}\MakeUppercase{Математика}}\par
\vspace{16mm}
{\Huge\bfseries\color{ink}Ревью домашней работы}\par
\vspace{11mm}
{\Large\color{accent}Кредиты, вклады и квадратные корни}\par
\vspace{24mm}
{\large Анастасия Павлова}\par
\vspace{6mm}
{\normalsize\color{muted}2 ноября 2026 года}\par
\end{center}
\vfill
\begin{center}
{\small Лёвин Артём Александрович}\par
{\small\color{muted}эксклюзивно для Анастасии Павловой}\par
\vspace{15mm}
\begin{tikzpicture}
\draw[accent!55,line width=.8pt] (-4.7,0) .. controls (-2.6,.7) and (-1.5,-.7) .. (0,0)
 .. controls (1.7,.55) and (3.2,-.6) .. (4.7,0);
\end{tikzpicture}
\end{center}
\clearpage
\section*{Сводная таблица}
\vspace{2mm}
В таблице перечислены только задания, в которых требуется исправить запись или решение. Нажмите на номер, чтобы перейти к подробному разбору.
\par\vspace{6mm}
\noindent\small
\renewcommand{\arraystretch}{1.45}
\begin{tabularx}{\linewidth}{@{}p{12mm}p{23mm}X p{25mm}p{25mm}@{}}
\toprule
\textbf{№} & \textbf{Тема} & \textbf{Суть затруднения} & \textbf{Ваша запись} & \textbf{Верная запись}\\
\midrule
\hyperref[task:7]{7} & Квадратный корень & В одном из промежуточных шагов корень приравнен к подкоренному числу; результат необходимо записать явно. & $\sqrt{1849}=1849$ & $\sqrt{1849}=43$\\
\addlinespace
\hyperref[task:10]{10} & Кредит & При выделении процентных начислений потеряно слагаемое, равное сумме кредита. & $p=50\%$ & $p=10\%$\\
\bottomrule
\end{tabularx}
\normalsize
\vfill
{\small\color{muted}В задании 7 правильная идея поиска числа по квадрату уже присутствует. В задании 10 необходимо внимательно отделять возврат основного долга от процентов.}
\clearpage
\thispagestyle{empty}
\vspace*{.36\textheight}
\begin{center}
{\Huge\bfseries\color{ink}Разбор ошибок}\par
\vspace{7mm}
{\normalsize\color{muted}Локализация, исправление, полное решение и одна задача на закрепление}
\end{center}
\clearpage
\phantomsection\label{task:7}
\section*{Задание 7. Вычисление квадратного корня}
\field{Текст задания}
Вычислите без калькулятора $\sqrt{5184}$, $\sqrt{1849}$ и $\sqrt{46656}$. Кратко объясните способ для каждого числа.
\field{Ваше решение и ответ}
Для первого числа использовано разложение на множители, получено $72$. Для третьего числа также использовано разложение и получено $216$. При работе со вторым числом встречается запись $\sqrt{1849}=1849$. Затем Вы находите границы между $40$ и $50$ и устанавливаете, что $43^2=1849$.
\field{Где именно возникает ошибка}
Последний достоверно правильный шаг для второго корня: рассматривается число $1849$ под знаком корня. Первый неверный переход:
\[
\textcolor{wrong}{\sqrt{1849}=1849}.
\]
\errorlabel\ Значение арифметического квадратного корня спутано с подкоренным выражением.
\field{Почему это неверно}
Запись $\sqrt{a}$ означает не само число $a$, а \emph{неотрицательное число, квадрат которого равен $a$} (при $a\ge 0$). Поэтому $\sqrt{1849}$ должно быть значительно меньше $1849$. Например, $\sqrt{100}=10$, поскольку $10^2=100$.

\keylabel\ Если Вы установили, что $43^2=1849$, то уже доказали, что $\sqrt{1849}=43$: число $43$ положительно и его квадрат равен $1849$.
\field{Исправление}
Вместо неверного равенства записываем цепочку:
\[
40^2=1600<1849<2500=50^2,\qquad 43^2=1849,
\]
\[
\boxed{\sqrt{1849}=43}.
\]
\field{Полное правильное решение}
\begin{align*}
5184&=72^2 &&\Longrightarrow\quad \sqrt{5184}=72,\\
1849&=43^2 &&\Longrightarrow\quad \sqrt{1849}=43,\\
46656&=216^2 &&\Longrightarrow\quad \sqrt{46656}=216.
\end{align*}
Последний корень удобно вычислить через разложение:
\[
46656=6^6,\qquad \sqrt{6^6}=6^3=216.
\]
\field{Проверка}
$72^2=5184$, $43^2=1849$, $216^2=46656$; все найденные значения неотрицательны.
\resultlabel\ Ответ: $72$, $43$, $216$.
\field{Задача на закрепление}
Вычислите без калькулятора $\sqrt{2209}$ и кратко поясните, как Вы нашли значение корня.
\clearpage
\phantomsection\label{task:10}
\section*{Задание 10. Дифференцированный кредит}
\field{Текст задания}
Дифференцированный кредит на $4$ года составил $1\,800\,000$ руб. Всего банку выплатили $2\,250\,000$ руб. Найдите годовую процентную ставку, первый и последний платежи.
\field{Ваше решение и ответ}
Вы верно начали с общей суммы выплат:
\[
2\,250\,000=1\,800\,000+1\,800\,000\cdot p\cdot\frac{4+1}{2},
\]
где $p$ — годовая ставка в долях единицы. Далее записано $2\,250\,000=4\,500\,000p$, откуда получена ставка $50\%$. С этой ставкой рассчитаны платежи $1\,350\,000$ руб. и $675\,000$ руб.
\field{Где именно возникает ошибка}
\errorlabel\ При переходе к уравнению для ставки исчез возврат основного долга:
\[
\underbrace{2\,250\,000=1\,800\,000+4\,500\,000p}_{\text{верно}}
\quad\not\Longrightarrow\quad
\underbrace{2\,250\,000=4\,500\,000p}_{\text{неверно}}.
\]
\field{Почему это неверно}
Сумма $2\,250\,000$ руб. включает две составляющие: возврат выданных $1\,800\,000$ руб. и проценты. Если сразу приравнять все выплаты к процентам, размер процентов окажется завышен. При переносе слагаемого $1\,800\,000$ в другую часть уравнения его нужно вычесть.
\keylabel\ Сначала выделяйте \emph{переплату}, затем находите процентную ставку. При равномерном погашении основного долга проценты каждого года начисляются на остаток перед соответствующим платежом.
\field{Исправление от последнего верного шага}
\begin{align*}
2\,250\,000&=1\,800\,000+4\,500\,000p,\\
2\,250\,000-1\,800\,000&=4\,500\,000p,\\
450\,000&=4\,500\,000p,\\
p&=\frac{450\,000}{4\,500\,000}=0{,}1=10\%.
\end{align*}
\field{Полное правильное решение}
За каждый из четырёх лет возвращается одинаковая часть основного долга:
\[
A=\frac{1\,800\,000}{4}=450\,000\ \text{руб.}
\]
В начале каждого года остаток долга равен соответственно
\[
1\,800\,000;\quad 1\,350\,000;\quad 900\,000;\quad 450\,000\ \text{руб.}
\]
Поэтому вся переплата по кредиту:
\begin{align*}
I&=2\,250\,000-1\,800\,000=450\,000\ \text{руб.},\\
I&=p(1\,800\,000+1\,350\,000+900\,000+450\,000)\\
 &=4\,500\,000p.
\end{align*}
Отсюда $p=0{,}1$, то есть $10\%$ годовых.

Первый платёж состоит из равной доли долга и процентов на всю первоначальную сумму:
\[
P_1=450\,000+0{,}1\cdot1\,800\,000=\boxed{630\,000\ \text{руб.}}
\]
Последний платёж содержит проценты только на остаток $450\,000$ руб.:
\[
P_4=450\,000+0{,}1\cdot450\,000=\boxed{495\,000\ \text{руб.}}
\]
\field{Проверка}
Промежуточные платежи: $P_2=585\,000$ руб., $P_3=540\,000$ руб. Их сумма вместе с первым и последним платежами:
\[
630\,000+585\,000+540\,000+495\,000=2\,250\,000.
\]
Получено точное совпадение с условием.
\resultlabel\ Годовая ставка — $10\%$; первый платёж — $630\,000$ руб.; последний — $495\,000$ руб.
\field{Задача на закрепление}
Кредит $1\,500\,000$ руб. выдан на $5$ лет с равномерным погашением основного долга. Всего банку выплатили $1\,950\,000$ руб. Найдите годовую процентную ставку, первый и последний платежи.
\end{document}